\documentclass[10pt,a4paper]{amsart}
\usepackage[margin=19mm]{geometry}
\usepackage{amsmath,amssymb,mathtools}
\usepackage[scr=rsfs,cal=euler]{mathalfa}
\usepackage{xfrac}
\usepackage[unicode,hidelinks,bookmarks=false]{hyperref}
\newcommand{\ie}{i.e.\ }
\newcommand{\cf}{cf.\ }
\newcommand{\resp}{respectively\ }
\newcommand{\Subsection}{\subsection}
\providecommand{\nobreakdash}{\nobreak\hbox{-}}
\setlength{\emergencystretch}{2em}
\multlinegap0pt
\usepackage{amssymb}
\usepackage{enumerate}
\usepackage[shortlabels]{enumitem}
\usepackage{mathtools}
\usepackage{float}
\usepackage{tikz}
\newtheorem{thm}{Theorem}
\newtheorem{lemma}{Lemma}
\newcommand{\N}{\mathbb{N}}
\newcommand{\Z}{\mathbb{Z}}
\title{On the Detection of Non-Roots of D'Arcais Polynomials}
\author{Bernhard Heim, Johann Stumpenhusen}
\date{Source: arXiv:2511.16276v1 (CC BY 4.0)}
\begin{document}
\maketitle

\noindent\emph{Source and licence.} On the Detection of Non-Roots of D'Arcais Polynomials. Version arXiv:2511.16276v1 (CC BY 4.0), distributed by arXiv under the Creative Commons Attribution 4.0 International licence, http://creativecommons.org/licenses/by/4.0/, as recorded in the arXiv licence metadata for this exact version and shown on the abstract page https://arxiv.org/abs/2511.16276v1. Full licence notice: \texttt{licenses/arxiv-2511.16276-CC-BY-4.0.txt}.\par\smallskip

\noindent\textit{Editorial scope.} Counted unit: the article's thm thm:NonRootsOnImaginaryAxisWithSigma (source0.tex lines 143--155). The article's complete original proof of it is delivered with it (source0.tex lines 502--519, one \texttt{proof} environment of 18 lines, which the article places away from the statement). No mathematics is written, repaired or re-derived: the statement and the proof are the article's own lines, in the article's own notation, and no retained line is substituted or re-flowed.  Retained uncounted as the source-local context the proof needs: lines 274--285; lines 441--454.  Nothing was omitted from any retained span, so the package carries no omission record.  No retained line contains a citation, so the wrapper carries no bibliography and no bibliography entry.  No retained line contains a figure environment, an image-inclusion command or drawing code, so no image is delivered and the package declares no asset dependency.  Editorial formatting only, in addition: an AMS \texttt{amsart} preamble that restates the article's own declarations and macros used by the retained lines, namely the environments \texttt{thm}, \texttt{lemma} on separate counters, together with the standard packages those lines need.  The retained environments print in this document as Theorem 1, Lemma 1 in order of appearance, whereas the article prints the same items with the numbering of its own full document.  The complete native source of the article as distributed in the arXiv e-print for this exact version is bundled unchanged as \texttt{sources/189505/source0.tex}.
\begin{thm}\label{thm:NonRootsTranslatingWithRationalIntegers}
    Let $a, b \in \Z$, $n \in \N$ and $g: \N \longrightarrow \Z$ be an arithmetic function such that $g(1) = 1$ and $m \geq 3$.
        \begin{enumerate}
            \item If $m$ is divisible by an odd rational prime and $2 \nmid a$, then
            $P_n^g(a\zeta_m + b) \neq 0.$
            \item If $m$ is divisble by a rational prime $p > 3$ or by $4$ and $3 \nmid a$ as well as $g(3) \equiv 0,1 \mod 3$, then
            $P_n^g(a\zeta_m + b) \neq 0.$
            \item Let $D \equiv 5 \mod 8$ be square-free. If $2 \nmid a$, then $P_n^g(a\omega_D + b) \neq 0.$
            \item Let $D \equiv 2 \mod 3$ be square-free. If $3 \nmid a$ and $g(3) \equiv 0, 1 \mod 3$, then
            $P_n^g(a\omega_D + b) \neq 0.$
        \end{enumerate}
\end{thm}

\begin{lemma}\label{lem:IndexOfaZeta_mAndaOmega_D}
    Let $a, b \in \Z$.
    \begin{enumerate}
        \item If $m \geq 3$, then
        \begin{equation*}
            \kappa_{a\zeta_m + b} = |a|^{\frac{\varphi(m)\left(\varphi(m) - 1\right)}{2}}
        \end{equation*}
        where $\varphi$ is the Euler totient function.
        \item If $D \in \Z$ square-free, then
        \begin{equation*}
            \kappa_{a\omega_D + b} = |a|. 
        \end{equation*}
    \end{enumerate}
\end{lemma}

\begin{thm}\label{thm:NonRootsOnImaginaryAxisWithSigma}
    Let $a, b \in \Z$ and $n \in \N$.
    \begin{enumerate}
        \item If $n \not \equiv 5 \mod 7$ and $21 \nmid a$, then $P_n^\sigma(ai + b) \neq 0$.
        \item If $n \equiv 5 \mod 7$ and either
        \begin{equation*}
            \textit{(i)} \; 3 \nmid a \qquad \textup{or} \qquad
            \textit{(ii)} \; a \not \equiv 0, \pm 1 \mod 7 \qquad \textup{or} \qquad
            \textit{(iii)} \; 7 \nmid a \textup{ and } 7 \nmid b,
        \end{equation*}
        then $P_n^\sigma(ai + b) \neq 0$.
    \end{enumerate}
\end{thm}

\begin{proof}[Proof of Theorem \ref{thm:NonRootsOnImaginaryAxisWithSigma}]
    The minimal polynomial of $ai + b$ is given by $m_{ai + b}(X) = (X - b)^2 + a^2$. Suppose $3 \nmid a$, then $P_n^\sigma(ai + b) \neq 0$ by virtue of the second item of Theorem \ref{thm:NonRootsTranslatingWithRationalIntegers} for all $n \in \N$. 
     
    Now assume $3 \mid a$. By our assumption on $a$, it follows $7 \nmid a$ and thus by Lemma \ref{lem:IndexOfaZeta_mAndaOmega_D}, we have $7 \nmid \kappa_{ai + b}$. Then $m_{ai + b}(X)$ does not decompose into irreducible linear polynomials modulo $7$ by the Dedekind--Kummer Theorem since $7 \equiv 3 \mod 4$ but instead remains a quadratic polynomial.
    \begin{enumerate}
        \item If $n \not \equiv 5, 6 \mod 7$, then $A_n^\sigma(X)$ decomposes into irreducible linear polynomials over $\Z$ and hence does so modulo $7$. Thus, $m_{ai + b}(X) \nmid P_n^\sigma(X)$.

        If $n \equiv 6 \mod 7$, then the only non-linear irreducible polynomial in the decomposition of $A_n^\sigma(X)$ modulo 7 is
        \[h_6(X) \equiv X^3 - X^2 - X + 4 \mod 7.\]
        Since $m_{ai +b}(X)$ and $h_6(X)$ are both irreducible, we again deduce $m_{ai+b}(X) \nmid P_n^\sigma(X)$.
        \item If $n \equiv 5 \mod 7$, the only non-linear irreducible polynomial in the decomposition of $A_n^\sigma(X)$ modulo 7 is
        \[h_5(X) \equiv X^2 + 1 \mod 7.\]
        If $7 \nmid b$, then $h_5(X)$ is obviously not associated to $m_{ai + b}(X)$ modulo 7 due to the latter having a non-vanishing linear term. If however $7 \mid b$, then
        \[m_{ai + b}(X) \equiv X^2 + a^2 \mod 7\]
        which is congruent to $h_5(X)$ if and only if $a^2 \equiv 1 \mod 7$.
    \end{enumerate}
    This completes the proof.
\end{proof}

\end{document}
